% Options for packages loaded elsewhere
\PassOptionsToPackage{unicode}{hyperref}
\PassOptionsToPackage{hyphens}{url}
\PassOptionsToPackage{dvipsnames,svgnames,x11names}{xcolor}
%
\documentclass[
  10pt,
  a4paper,
  DIV=11]{scrartcl}

\usepackage{amsmath,amssymb}
\usepackage{iftex}
\ifPDFTeX
  \usepackage[T1]{fontenc}
  \usepackage[utf8]{inputenc}
  \usepackage{textcomp} % provide euro and other symbols
\else % if luatex or xetex
  \usepackage{unicode-math}
  \defaultfontfeatures{Scale=MatchLowercase}
  \defaultfontfeatures[\rmfamily]{Ligatures=TeX,Scale=1}
\fi
\usepackage[]{libertine}
\ifPDFTeX\else  
    % xetex/luatex font selection
    \setmainfont[]{Arial}
\fi
% Use upquote if available, for straight quotes in verbatim environments
\IfFileExists{upquote.sty}{\usepackage{upquote}}{}
\IfFileExists{microtype.sty}{% use microtype if available
  \usepackage[]{microtype}
  \UseMicrotypeSet[protrusion]{basicmath} % disable protrusion for tt fonts
}{}
\makeatletter
\@ifundefined{KOMAClassName}{% if non-KOMA class
  \IfFileExists{parskip.sty}{%
    \usepackage{parskip}
  }{% else
    \setlength{\parindent}{0pt}
    \setlength{\parskip}{6pt plus 2pt minus 1pt}}
}{% if KOMA class
  \KOMAoptions{parskip=half}}
\makeatother
\usepackage{xcolor}
\usepackage[top=30mm,left=20mm,heightrounded]{geometry}
\setlength{\emergencystretch}{3em} % prevent overfull lines
\setcounter{secnumdepth}{3}
% Make \paragraph and \subparagraph free-standing
\makeatletter
\ifx\paragraph\undefined\else
  \let\oldparagraph\paragraph
  \renewcommand{\paragraph}{
    \@ifstar
      \xxxParagraphStar
      \xxxParagraphNoStar
  }
  \newcommand{\xxxParagraphStar}[1]{\oldparagraph*{#1}\mbox{}}
  \newcommand{\xxxParagraphNoStar}[1]{\oldparagraph{#1}\mbox{}}
\fi
\ifx\subparagraph\undefined\else
  \let\oldsubparagraph\subparagraph
  \renewcommand{\subparagraph}{
    \@ifstar
      \xxxSubParagraphStar
      \xxxSubParagraphNoStar
  }
  \newcommand{\xxxSubParagraphStar}[1]{\oldsubparagraph*{#1}\mbox{}}
  \newcommand{\xxxSubParagraphNoStar}[1]{\oldsubparagraph{#1}\mbox{}}
\fi
\makeatother


\providecommand{\tightlist}{%
  \setlength{\itemsep}{0pt}\setlength{\parskip}{0pt}}\usepackage{longtable,booktabs,array}
\usepackage{calc} % for calculating minipage widths
% Correct order of tables after \paragraph or \subparagraph
\usepackage{etoolbox}
\makeatletter
\patchcmd\longtable{\par}{\if@noskipsec\mbox{}\fi\par}{}{}
\makeatother
% Allow footnotes in longtable head/foot
\IfFileExists{footnotehyper.sty}{\usepackage{footnotehyper}}{\usepackage{footnote}}
\makesavenoteenv{longtable}
\usepackage{graphicx}
\makeatletter
\def\maxwidth{\ifdim\Gin@nat@width>\linewidth\linewidth\else\Gin@nat@width\fi}
\def\maxheight{\ifdim\Gin@nat@height>\textheight\textheight\else\Gin@nat@height\fi}
\makeatother
% Scale images if necessary, so that they will not overflow the page
% margins by default, and it is still possible to overwrite the defaults
% using explicit options in \includegraphics[width, height, ...]{}
\setkeys{Gin}{width=\maxwidth,height=\maxheight,keepaspectratio}
% Set default figure placement to htbp
\makeatletter
\def\fps@figure{htbp}
\makeatother

%% load packages
\usepackage{sectsty}
\usepackage{fontspec}
\usepackage{scrlayer-scrpage}
\usepackage{fontspec}

% Define title formatting

%% Set font sizes for sections and subsections
\sectionfont{\fontsize{11}{16}\selectfont\bfseries} % Arial 14pt bold
\subsectionfont{\fontsize{10}{14}\selectfont} % Arial 12pt

\lehead[test1]{\pagemark}
\cehead[]{}
\rehead[test 2]{\leftmark}
\lohead[{\fontsize{8}{10}\selectfont\upshape\textbf{Institut Sekundarstufe I}} \\
\fontsize{8}{10}\selectfont\upshape Fabrikstrasse 8, CH-3012 Bern \\
\fontsize{8}{10}\selectfont\upshape T +41 31 309 21 15, contactdesk@phbern.ch, www.phbern.ch
]{\rightmark}
\rohead[PHBern]{}
\KOMAoption{captions}{tableheading}
\makeatletter
\@ifpackageloaded{caption}{}{\usepackage{caption}}
\AtBeginDocument{%
\ifdefined\contentsname
  \renewcommand*\contentsname{Inhaltsverzeichnis}
\else
  \newcommand\contentsname{Inhaltsverzeichnis}
\fi
\ifdefined\listfigurename
  \renewcommand*\listfigurename{Abbildungsverzeichnis}
\else
  \newcommand\listfigurename{Abbildungsverzeichnis}
\fi
\ifdefined\listtablename
  \renewcommand*\listtablename{Tabellenverzeichnis}
\else
  \newcommand\listtablename{Tabellenverzeichnis}
\fi
\ifdefined\figurename
  \renewcommand*\figurename{Abbildung}
\else
  \newcommand\figurename{Abbildung}
\fi
\ifdefined\tablename
  \renewcommand*\tablename{Tabelle}
\else
  \newcommand\tablename{Tabelle}
\fi
}
\@ifpackageloaded{float}{}{\usepackage{float}}
\floatstyle{ruled}
\@ifundefined{c@chapter}{\newfloat{codelisting}{h}{lop}}{\newfloat{codelisting}{h}{lop}[chapter]}
\floatname{codelisting}{Listing}
\newcommand*\listoflistings{\listof{codelisting}{Listingverzeichnis}}
\makeatother
\makeatletter
\makeatother
\makeatletter
\@ifpackageloaded{caption}{}{\usepackage{caption}}
\@ifpackageloaded{subcaption}{}{\usepackage{subcaption}}
\makeatother

\ifLuaTeX
\usepackage[bidi=basic]{babel}
\else
\usepackage[bidi=default]{babel}
\fi
\babelprovide[main,import]{ngerman}
\ifPDFTeX
\else
\babelfont{rm}[]{Arial}
\fi
% get rid of language-specific shorthands (see #6817):
\let\LanguageShortHands\languageshorthands
\def\languageshorthands#1{}
\ifLuaTeX
  \usepackage{selnolig}  % disable illegal ligatures
\fi
\usepackage{bookmark}

\IfFileExists{xurl.sty}{\usepackage{xurl}}{} % add URL line breaks if available
\urlstyle{same} % disable monospaced font for URLs
\hypersetup{
  pdftitle={Gleichungen},
  pdfauthor={Richard Conrardy},
  pdflang={de},
  colorlinks=true,
  linkcolor={(3,1.5)},
  filecolor={Maroon},
  citecolor={Blue},
  urlcolor={Blue},
  pdfcreator={LaTeX via pandoc}}


\title{Gleichungen}
\usepackage{etoolbox}
\makeatletter
\providecommand{\subtitle}[1]{% add subtitle to \maketitle
  \apptocmd{\@title}{\par {\large #1 \par}}{}{}
}
\makeatother
\subtitle{Einführung und Theorie}
\author{David Grossenbacher}
\date{22.09.2024}

\begin{document}
\maketitle
\begin{abstract}
Eine Doppellektion zur Einführung in Gleichungen
\end{abstract}

\renewcommand*\contentsname{Inhaltsverzeichnis}
{
\hypersetup{linkcolor=}
\setcounter{tocdepth}{3}
\tableofcontents
}

\section{}\label{section}

\section{Lehrerkommentar}\label{lehrerkommentar}

\href{https://be.lehrplan.ch/101E8xmhHHgram3LHqgCx2Gwdg3vJ6cPY}{MA.1.A.4.i}
können ein Produkt mit gleichen Faktoren als Potenz schreiben und
umgekehrt (z.B. \(15 \cdot 15 \cdot 15 = 15^3\) ;
\(a \cdot a \cdot a \cdot a = a^4\) ).

\href{https://be.lehrplan.ch/101AJt2yMqVxRygKJ2xZqrbaqzAhfYnDL}{MA.1.A.4.j}
Erweiterung: können lineare Gleichungen mit einer Variablen mit
Äquivalenzumformungen lösen (z.B. \(5x + 3 = 7\))
\href{https://be.lehrplan.ch/101eyUHN5UTpNM7eNmb25PVMKdrYH99AJ}{MA.3.c.3.g}
Erweiterung: können Buchstabenterme, Formeln und lineare
Funktionsgleichungen mit Sachsituationen konkretisieren (z.B. die
Funktionsgleichung y = 2x + 3 mit Preis = 2 · Anzahl + 3).

\section{Lehrpersonenkommentar}\label{lehrpersonenkommentar}

Die vorliegende Lernumgebung dient als Einführung und erste Annäherung
an das Thema «Gleichungen» für 7. Klässler im ersten Semester. Die
beigelegte Lektionsplanung ist für etwa 3-4 Lektionen ausgelegt.

\subsection{Vorangehende Kompetenzen}\label{vorangehende-kompetenzen}

Folgende \textbf{vorangehende Kompetenzen} sollten die SuS behandelt
haben:

\begin{itemize}
\tightlist
\item
  Die SuS wissen, was eine Variable ist und sie kennen die einfachen
  Term-Vereinfachungsregeln (Bspw.: \(x + 2x = 3x\)).
\item
  Die SuS beherrschen die Grundoperationen.
\end{itemize}

\subsection{Grundlegende
Stolpersteine}\label{grundlegende-stolpersteine}

Weiter sind die folgenden grundlegenden \textbf{Stolpersteine} zu
beachten:

\begin{itemize}
\tightlist
\item
  Die SuS kennen die Rechenregeln von ganzen und rationalen Zahlen noch
  nicht. Die Lösungen und die in den Aufgaben gebrauchten Zahlen sind
  also immer rein natürliche Zahlen. Folgendes Beispiel ist
  dementsprechend nicht Inhalt der LU:
\end{itemize}

\begin{figure}[H]

{\centering \includegraphics{.pdf}

}

\caption{Bild}

\end{figure}%

\begin{itemize}
\tightlist
\item
  Die SuS können Gleichungen noch nicht umstellen. Die Variable \(x\)
  wird durch reines Einsetzen und Kopfrechnen ermittelt.
\item
  Die SuS haben noch kein oder ein falsches Verständnis dafür, was das
  Gleichheitszeichen bedeutet. In der Primarstufe wird oftmals gesagt
  ``eins und eins gibt zwei''. Daraus könnte sich das Fehlkonzept
  entwickeln, dass \(1 + 1\) irgendwie zu \(2\) führt. Das Verständnis
  dafür, dass beide Terme auf der linken und rechten Seite des
  Gleichheitszeichens ausgewertet gleich groß sein müssen, wäre
  dementsprechend nicht vorhanden.
\item
  Termumformungen sind in dieser LU kein Thema.
\end{itemize}

\subsection{Lernziele}\label{lernziele}

In dieser Lernumgebung werden folgende Lernziele angepeilt:

\begin{itemize}
\item
  Die SuS können einfache Gleichungen mit einer Variable durch reines
  Einsetzen und Kopfrechnen ermitteln.
\item
  Die SuS können Gleichungen ausgehend von Texten aufschreiben.
\item
  Die SuS füllen Wertetabellen für Gleichungen mit 2 Variablen korrekt
  aus.
\item
  Die Schüler erstellen keine ``Gleichheitsketten'' und verwenden das
  Äquivalenzzeichen korrekt. Bsp:

  \(3x + 1 = 10 \leftrightarrow x = 3\)

  und nicht

  \(3x + 1 = 10 = x = 3\)
\end{itemize}

\subsection{Arbeitsunterlagen für
SuS}\label{arbeitsunterlagen-fuxfcr-sus}

\begin{itemize}
\tightlist
\item
  Da das Thema ``Gleichungen'' erst eingeführt wird, haben die Übungen
  nicht sehr viel Anwendungscharakter.
\item
  Die Aufgabenstellungen sind nicht sehr differenziert, sie sollten aber
  in allen Klassen für die Erarbeitung der Grundlagen hilfreich sein.
  Bei sehr heterogenen Klassen müssen allenfalls zusätzliche
  Arbeitsunterlagen den Bedürfnissen entsprechend angeschafft werden.
\item
  Das AB1 ist relativ repetitiv. Es muss nicht das gesamte AB gelöst
  werden, bevor die SuS zu Übung 2 übergehen.
\item
  Die Übung 3 ist grundsätzlich nach der Einführung zu den Wertetabellen
  zu bearbeiten. Nach eigenem Ermessen dürfen die SuS auch noch an den
  vorherigen Aufträgen weitermachen.
\item
  Die unten aufgeführten .html Dateien orientieren sich an der
  Aufgabenstellung des AB1 (\(x\) herausfinden) und an den ersten 5
  Wertetabellen-Aufgaben (Wertetabellen\_Generator).
\end{itemize}

\subsubsection{Zusatzaufgaben für schnelle
SuS:}\label{zusatzaufgaben-fuxfcr-schnelle-sus}

\begin{itemize}
\tightlist
\item
  Erstelle eigene Textaufgaben (siehe Übung 2).
\item
  Erstelle einen Merkhefteintrag, wo die wichtigsten theoretischen
  Inhalte festgehalten werden.
\item
  Übe weiter in den Rechenprogrammen.
\end{itemize}

\subsubsection{Weitere Übungsmöglichkeiten, sollte die gesamte Klasse
sehr schnell
arbeiten:}\label{weitere-uxfcbungsmuxf6glichkeiten-sollte-die-gesamte-klasse-sehr-schnell-arbeiten}

\begin{itemize}
\tightlist
\item
  Die von den SuS erstellten Textaufgabenbeispiele können im Plenum
  zusammen gelöst werden.
\end{itemize}

\subsection{Unterrichtsmaterialien}\label{unterrichtsmaterialien}

Die LP benötigt eine Visualisierungshilfe (Visualiser, Tafel, o.ä.). Die
Arbeitsaufgaben für die SuS können in dieser Form weiterverwendet und
ausgedruckt werden. Weiter sollen A5-Blätter für die Zusatzaufgabe in
Übung 2 bereitgehalten werden.

\subsection{Ergebnissicherung}\label{ergebnissicherung}

Es ist keine konkrete Art der Ergebnissicherung eingeplant. Es ist
Auftrag der LP, während der individuellen Lernphasen Erkenntnisse über
den Lernstand der individuellen SuS zu sammeln. Für die Lernprozess- /
formative Beurteilung kann ein Merkhefteintrag eingefordert werden.

\subsection{Unterrichtsmethoden}\label{unterrichtsmethoden}

Makromethodisch findet der Unterricht in einem fremdgesteuerten Setting
statt. Die LP instruiert und die SuS führen vorbereitete
Aufgaben/Übungen aus.

\begin{itemize}
\tightlist
\item
  \textbf{Direkte Instruktionen:}

  \begin{itemize}
  \tightlist
  \item
    Es sind zwei Frontalunterrichts-Sequenzen eingeplant. Hier werden
    hauptsächlich theoretische Inhalte vermittelt und auch kurze
    Einstiegsbeispiele gezeigt.
  \item
    Es wird versucht, am Vorverständnis der SuS anzuknüpfen.
  \item
    Die Bedeutung des Gleichheitszeichens und Äquivalenzpfeils werden im
    Plenum vermittelt.
  \end{itemize}
\item
  \textbf{Partnerarbeiten:}

  \begin{itemize}
  \tightlist
  \item
    Während den Frontalunterrichtssequenzen werden kurze Murmelphasen
    eingeplant, um die Anteilnahme der SuS zu fördern. Hier dürfen sich
    die SuS kurz miteinander austauschen.
  \item
    Während den Arbeitsphasen können sich die SuS zusammensetzen und
    Probleme gegebenenfalls miteinander betrachten.
  \end{itemize}
\item
  \textbf{Individuelles Lernen (Offener Lernunterricht):}

  \begin{itemize}
  \tightlist
  \item
    Die SuS arbeiten selbstständig an den Übungsaufgaben. Hier können
    sie sich im eigenen Tempo in das Gehörte vertiefen.
  \item
    Die LP agiert als Coach und hilft wo nötig.
  \end{itemize}
\item
  \textbf{Digitale Medien:}

  \begin{itemize}
  \tightlist
  \item
    Die Aufgaben können den SuS digital gezeigt werden.
  \item
    Es gibt zwei verschiedene Rechenprogramme, mit Hilfe von welchen SuS
    weiterüben können.
  \end{itemize}
\item
  \textbf{Deduktive Vorgehensweise:}

  \begin{itemize}
  \tightlist
  \item
    Es werden zuerst theoretische Inhalte vermittelt. In den
    Übungsaufgaben werden dann eher konkrete Fälle betrachtet.
  \end{itemize}
\end{itemize}

\subsection{Unterrichtstechniken}\label{unterrichtstechniken}

In der vorliegenden Unterrichtsplanung werden verschiedene
Unterrichtstechniken eingesetzt.

\begin{itemize}
\tightlist
\item
  \textbf{Visualisierungshilfen:}

  \begin{itemize}
  \tightlist
  \item
    Um das Erzählte zu unterstützen, zeigt die LP an der Wandtafel (oder
    am Visualiser, o.ä.) die Rechnungen und Skizzen auf.
  \end{itemize}
\item
  \textbf{Fragetechniken:}

  \begin{itemize}
  \tightlist
  \item
    Während den Einführungen werden verschiedene Fragen gestellt:

    \begin{itemize}
    \tightlist
    \item
      \emph{Rhetorische Fragen:} Es werden Fragen aufgeworfen, welche
      noch keine Antwort von den SuS verlangen. Bsp.: ``Was bedeutet das
      Gleichheitszeichen?'' Eine konkrete Antwort wird erst nach dem
      Betrachten der Waage angeschaut.
    \item
      \emph{Konkrete Fragen:} Es werden im Plenum Fragen gestellt,
      welche dann im Plenum besprochen werden.
    \end{itemize}
  \end{itemize}
\item
  \textbf{Aktivierungstechniken:}

  \begin{itemize}
  \tightlist
  \item
    Um Antworten bei den jeweiligen Fragen zu generieren, werden
    hauptsächlich Murmelgruppen verwendet. Hier besprechen die SuS mit
    ihrem Pultnachbar die Fragen, um den Einstieg zu erleichtern.
  \end{itemize}
\item
  \textbf{Reflexionstechniken:}

  \begin{itemize}
  \tightlist
  \item
    Es besteht die Möglichkeit, dass die SuS einen Merkhefteintrag
    erstellen, in welchem sie das Gelernte nochmals festhalten.
  \end{itemize}
\item
  \textbf{Schritt-für-Schritt Erklärungen:}

  \begin{itemize}
  \tightlist
  \item
    Die SuS werden in verschiedenen Aufgaben angeleitet, damit sie das
    Lösungsvorgehen bei Gleichungen verstehen lernen.
  \end{itemize}
\end{itemize}

\subsection{Didaktische Prinzipien
(Mathematik)}\label{didaktische-prinzipien-mathematik}

\begin{itemize}
\tightlist
\item
  \textbf{Anknüpfen an Vorwissen + Einführung in neue Konzepte}

  \begin{itemize}
  \tightlist
  \item
    Die Lernphase (LP) beginnt mit der Wiederholung des bereits
    bekannten Themas „Variablen``.
  \item
    Dies stellt sicher, dass das Vorwissen aktiviert wird und die
    Schülerinnen und Schüler (SuS) die Möglichkeit haben, auf bereits
    Gelerntes zurückzugreifen.
  \item
    Danach wird schrittweise das neue Konzept „Gleichungen`` eingeführt.
  \end{itemize}
\item
  \textbf{Verständnis von Begriffen und Symbolen (Variable, Term,
  Gleichheitszeichen, etc.)}

  \begin{itemize}
  \tightlist
  \item
    Bezüglich der Sprachsensibilität im Mathematikunterricht müssen
    diese Begriffe im Unterricht explizit verwendet werden.
  \item
    Dies fördert ein klares Verständnis und eine präzise Verwendung von
    mathematischen Begriffen.
  \end{itemize}
\item
  \textbf{Förderung des funktionalen Denkens}

  \begin{itemize}
  \tightlist
  \item
    Die LP zeigt beispielsweise lineare Beziehungen auf.
  \item
    Es wird erklärt, wie sich die Werte in der Zweierreihe durch
    Multiplikation von (x) mit 2 ergeben.
  \item
    Dieses Verständnis legt die Basis für weiterführende Themen wie
    lineare Funktionen und Graphen.
  \end{itemize}
\item
  \textbf{Veranschaulichung der algebraischen Strukturen}

  \begin{itemize}
  \tightlist
  \item
    Durch das Hinzufügen weiterer Summanden zu einer Gleichung (von
    (y=2x) zu (y=2x+3)) wird den SuS verdeutlicht, wie sich die
    mathematische Struktur verändert.
  \item
    Dies zeigt auf, was diese Veränderung für die Lösungsmenge bedeutet.
  \item
    Dies ist ein wichtiger Aspekt der algebraischen Denkweise, die es
    den SuS ermöglicht, allgemeine Strukturen in Gleichungen zu erkennen
    und zu analysieren.
  \end{itemize}
\item
  \textbf{Exploratives Arbeiten}

  \begin{itemize}
  \tightlist
  \item
    Um der rein deduktiven Vorgehensweise entgegenzuwirken, können die
    SuS eigene Textaufgabenbeispiele erstellen.
  \item
    Hier geht es darum, Alltagsbeispiele zu mathematisieren.
  \end{itemize}
\end{itemize}




\end{document}
